\section{Series Arc-Fault Detection}
\label{sec:ch4-series-arc-detection}

The series arc-fault tests evaluated whether waveform disturbance could be
converted into a relay decision. During arcing, the current waveform was
expected to show low-current gaps, restrike spikes, unstable cycle shape, and
changes in mid/high-frequency content. These behaviors correspond to the
features defined in Section~\ref{sec:ch3-waveform-features}, especially
maximum low-current run, low-current ratio, zero-dwell behavior, midband
residual ratio, and spectral flux.

The trial-level results are summarized as confusion matrices in
Table~\ref{tab:ch4-arc-confusion}. Here, ``relay trips'' is the positive
prototype response and ``arcing event'' is the positive ground-truth condition.
\begin{table}[H]
  \centering
  \caption{Trial-level confusion matrices for series arc-fault testing}
  \label{tab:ch4-arc-confusion}
  \begin{tabular}{p{0.26\textwidth}p{0.24\textwidth}cc}
    \toprule
    \textbf{Device group} & \textbf{Prototype response} & \textbf{Arcing event} & \textbf{Normal event} \\
    \midrule
    Recognized loads & Relay trips & 49 & 0 \\
    Recognized loads & Relay remains on & 1 & 50 \\
    \midrule
    Unrecognized loads & Relay trips & 47 & 1 \\
    Unrecognized loads & Relay remains on & 3 & 49 \\
    \midrule
    All tested loads & Relay trips & 96 & 1 \\
    All tested loads & Relay remains on & 4 & 99 \\
    \bottomrule
  \end{tabular}
\end{table}

Table~\ref{tab:ch4-arc-confusion} shows that the prototype detected 96 of 100
arc trials and produced one relay trip during 100 normal trials. The recognized
load group had no normal-event false trip and one missed arc event. The
unrecognized group was more difficult, with three missed arcs and one false
trip. This difference is expected because unrecognized devices introduce
waveform behavior not directly represented during model fitting.

The result supports the arc-detection objective under the tested laboratory
conditions. The classifier response was strongest when the arc created
persistent low-current gaps and spectral disturbance, allowing the leaky arc
counter to reach the relay-trip threshold. The missed trials indicate that not
all arcs produced equally strong or persistent feature responses, especially
when the load type differed from the fitted set. Therefore, the prototype
should be interpreted as a controlled feasibility demonstration, not as a
certified AFCI-equivalent device.

\section{Arc-Fault Dataset Summary}
\label{sec:ch4-arc-dataset-summary}

The cleaned arc-training dataset contained 49,296 rows, of which 283 were arc
markers, equivalent to about 0.577\% of the retained data. This confirms that
the problem was highly imbalanced: most frames represented normal operation
even in sessions that contained an arc event. The imbalance explains why
frame-level accuracy alone can be misleading; a classifier can appear accurate
by predicting mostly normal frames. Precision, recall, F1-score, and
event-level metrics are therefore needed to interpret the safety behavior.

Complete trial tables and raw collection logs are provided in
Appendices~\ref{app:e} and~\ref{app:f}. These appendices are important because
the summary table in the main text cannot show every feature frame, arc marker,
or relay-state transition.
